\section{De sequential data a sequential decision making}

Antes de entrar en las fórmulas de RL, conviene entender por qué el ponente empieza hablando de lenguaje, visión, proteínas y generación de código. La experiencia común de estos campos no es que «el Transformer lo puede todo», sino que el modelado de secuencias supervisado o autosupervisado a gran escala suele tener una loss relativamente suave, técnicas de regularización maduras e infraestructura de entrenamiento reutilizable. Lo que el ponente quiere trasladar realmente a los problemas de decisión es esa interfaz de entrenamiento y esa ruta de escalamiento.

\subsection{La abstracción común detrás del éxito del Transformer}

El lenguaje, los patch de imagen, los residuos de proteínas y los token de código parecen distintos, pero todos pueden representarse como secuencias y aprender distribuciones condicionales mediante masked prediction o autoregressive prediction. Al leer la siguiente slide, no compare primero el número de parámetros de cada modelo; observe antes cómo convierten los objetos de cada dominio en token y cómo los modelan con la misma clase de next-token machinery.

\lecturefigure{slide-02-transformers-for-sequential-data.jpg}{El Transformer ya se ha usado con muchos tipos de sequential data, como lenguaje, visión, proteínas y código.}{Video oficial de Stanford Online, 01:40.}

\begin{knowledgebox}{Lectura de la figura: lo común no es la modalidad de entrada, sino la interfaz de entrenamiento}
GPT-3, ViT, AlphaFold y los modelos de código de la figura no comparten la misma semántica de entrada, pero todos codifican objetos complejos como representaciones componibles y aprenden relaciones condicionales mediante optimización por gradiente con lotes grandes. Con este conjunto de ejemplos, el ponente plantea una pregunta de ingeniería: si las trayectorias de decisión también pueden convertirse en secuencias, ¿se pueden reutilizar objetivos de entrenamiento, optimizadores y scaling recipe más estables? Se trata de una hipótesis de investigación, no de un teorema que se derive automáticamente de estos casos de éxito.
\end{knowledgebox}

\subsection{Por qué importa escalar de forma estable}

La ventaja de los modelos a gran escala no es solo que tengan «más parámetros». Lo decisivo es si, al aumentar los datos y el presupuesto de cómputo, el comportamiento del entrenamiento es predecible, si las fallas pueden localizarse, si la regularización es madura y si la loss puede bajar de forma estable sin necesidad de un bootstrap complejo. La dificultad de RL muchas veces no es la falta de capacidad expresiva, sino que el objetivo, la distribución de los datos y los cambios en la policy están acoplados entre sí, de modo que el sistema de entrenamiento persigue continuamente la nueva distribución que él mismo genera.

\lecturefigure{slide-03-scalable-models-stable-training.jpg}{Las curvas de entrenamiento de los sequence model a gran escala suelen ser más suaves y, en términos de ingeniería, más fáciles de escalar.}{Video oficial de Stanford Online, 02:20; la slide cita a Kaplan et al. 2020.}

\teachervoice{El ponente enfatiza que las stable training dynamics son importantes para el practitioner: basta con tener más datos y más cómputo para ampliar el modelo paso a paso y observar una mejora relativamente suave de la loss. Una de las motivaciones de Decision Transformer es llevar esa «capacidad de entrenamiento» al decision making, y no solo buscar un nuevo módulo de red.}

\begin{warningbox}{Una loss estable no equivale a una calidad de decisión estable}
Que la pérdida supervisada baje solo indica que el modelo ajusta mejor la distribución condicional de acciones dentro de la distribución de entrenamiento. El retorno en el entorno real depende además de la distribución de estados que produce el rollout, de la acumulación de errores, de si el retorno objetivo es alcanzable y de la aleatoriedad del entorno. Más adelante habrá que medir por separado el realized environmental return; la loss de entrenamiento no puede sustituir a la policy evaluation.
\end{warningbox}

\subsection{De la percepción a la acción}

Los modelos de percepción producen etiquetas, texto o representaciones; la salida de un modelo de decisión modifica el entorno y determina qué se verá en el siguiente paso. Este ciclo cerrado hace que el sequential decision making tenga, a diferencia de la predicción estática, distribution shift y credit assignment. Cuando el ponente señala el robot con «un cilindro cargado de conocimiento visual y lingüístico», no quiere decir que al robot solo le falte un Transformer, sino subrayar que el perception knowledge debe convertirse finalmente en una action sequence.

\lecturefigure{slide-04-sequential-data-to-sequential-decision.jpg}{La pregunta de investigación pasa de sequential data a sequential decision making.}{Video oficial de Stanford Online, 03:10.}

\begin{importantbox}{El giro central}
El modelado estático de secuencias aprende $p(x_t\mid x_{<t})$; el modelado de secuencias de decisión necesita aprender
\[
p(a_t\mid \text{goal},s_{\le t},a_{<t},r_{<t}),
\]
donde la acción $a_t$ modifica el estado futuro $s_{t+1}$. Por eso, que «la tokenización sea similar» no elimina los riesgos del closed-loop control; solo ofrece una interfaz unificada de aproximación de funciones y de entrenamiento.
\end{importantbox}

\subsection{Resumen de la sección}

El punto de partida de Decision Transformer es trasladar un paradigma de entrenamiento: convertir las trayectorias de RL en secuencias con la esperanza de reutilizar los objetivos estables y la capacidad de escalamiento del sequence modeling. Sin embargo, la decisión tiene retroalimentación en ciclo cerrado; el éxito en dominios estáticos no basta para demostrar que el método funciona, y es necesario pasar al MDP, a los datos fuera de línea y a la evaluación con rollout reales.
